\section{Thermoelastic damping in flexure: Zener and exact Fourier solutions}\label{sec:flexted}

Longitudinal waves (\cref{sec:classical}) are only one of the settings in which the thermoelastic loss appears. In micro-resonators the dominant case is flexure of a thin beam, where the temperature gradient across the thickness $h$ relaxes by heat conduction in the transverse direction.

\subsection{Two classical results}
Zener's single-relaxation treatment gives
\begin{equation}\label{eq:zenerQ}
Q^{-1}=\Delta\,\frac{\omega\tau_Z}{1+(\omega\tau_Z)^2},\qquad \tau_Z=\frac{h^2}{\pi^2\chi},\qquad \Delta=\frac{E\alpha_T^2T_0}{c_v},
\end{equation}
with $\chi$ the thermal diffusivity, $\alpha_T$ the linear expansion coefficient and $c_v$ the volumetric heat capacity. The exact solution of the Fourier problem for a thin rectangular beam in Euler--Bernoulli flexure, due to Lifshitz and Roukes, is
\begin{equation}\label{eq:LR}
Q^{-1}=\Delta\left[\frac{6}{\xi^{2}}-\frac{6}{\xi^{3}}\,\frac{\sinh\xi+\sin\xi}{\cosh\xi+\cos\xi}\right],\qquad \xi=h\sqrt{\frac{\omega}{2\chi}} .
\end{equation}
Both reduce to $Q^{-1}\simeq\Delta\,\omega\tau_Z$ for $\omega\tau_Z\ll1$.

\subsection{Quantitative comparison}
\Cref{fig:flexted} compares \cref{eq:zenerQ,eq:LR} (script \path{scripts/fig_flexTED.py}). With $x=\omega\tau_Z$ the exact solution peaks at $Q^{-1}/\Delta=0.494$ near $x=1.008$, against $0.5$ at $x=1$ for the Zener form. The relative error of the Zener form is $+1.3\%$ at $x=0.01$, $0.1$ and $1$, $-5\%$ at $x=10$ and $-14\%$ at $x=100$. In practice the single-relaxation form is therefore accurate to a few percent through the peak and on its low-frequency side, and underestimates the loss at $x\gtrsim10$, where the exact solution decays as $x^{-1/2}$ whereas the Zener form decays as $x^{-1}$.

\begin{figure}[t]
\centering
\includegraphics[width=\linewidth]{figs/fig_flexTED.pdf}
\caption{(a) Thermoelastic loss of a thin beam in flexure: exact Fourier solution \cref{eq:LR} and the Zener form \cref{eq:zenerQ}. (b) Relative error of the Zener form.}
\label{fig:flexted}
\end{figure}

\subsection{Where hyperbolic corrections enter}
\Cref{eq:LR} assumes parabolic heat conduction across the thickness. Replacing the Fourier law by a lagging law (\cref{sec:lags}) changes the result only when the transverse relaxation rate $\pi^2\chi/h^2$ approaches the inverse lag. For $\tau_q\sim10^{-12}$\,s and $h\sim100$\,nm, with $\chi\sim10^{-4}$\,m$^2$/s, the product $\omega\tau_q$ at the Zener peak is of order $10^{-2}$, so the parabolic result is adequate; corrections matter for $h$ well below $\sim 10$\,nm or when ballistic transport is relevant, in which case the Guyer--Krumhansl form (\cref{sec:kinetic}) is the appropriate next step rather than a pure lag law.
